\section{Method}\label{sec:method}
This is a critical (theoretical) review, not a PRISMA-complete systematic review. It selects the literatures needed to evaluate one mechanistic argument and states what each can and cannot show. Searches ran in October 2026 across PubMed, Europe PMC, Crossref, publisher pages and the reference lists of key meta-analyses, in eight slices, the first four built for the executive-function argument: (a) efficacy, bias and head-to-head comparisons; (b) mechanisms and component (dismantling) studies, including exposure and third-wave alternatives; (c) executive function, stress and emotion regulation, including executive function as a moderator of treatment; and (d) ADHD (CBT, medication, mindfulness, executive function, anxiety comorbidity); and, added in Draft 3 to cover all the reasons CBT is said to fail, (e) non-response, dropout, relapse and harm; (f) patient and context factors (comorbidity, neurodivergence, culture, access); (g) therapist, delivery and real-world effectiveness; and (h) theory, mechanism and measurement critiques. Priority was given to meta-analyses, umbrella reviews, component trials and moderator studies, and to the strongest critiques and defences of CBT.

\paragraph{Verification} Every reference was resolved through Crossref (title, authors and year compared with the intended citation). Every quoted number was read from the abstract or full text, not recalled. Roughly 440 candidate sources were verified this way; several candidate references proposed from memory were dropped because the DOI did not resolve or resolved to a different paper. Sources for which only a partial abstract or a secondary summary was available are flagged in \texttt{deep-research/sources/}. Popular sources (books, talks, blogs) are cited only for the claims their authors make.

\paragraph{Limitations of the method} Single-reviewer screening; no pre-registered protocol; no formal risk-of-bias assessment of included reviews; evidence from depression is used where anxiety evidence on a question is thin, and flagged when so. A PRISMA-style umbrella search is the obvious next step (Section~\ref{sec:agenda}).
